\title{Byte Latent Transformer: Patches Scale Better Than Tokens}

\begin{abstract}
We present the Byte Latent Transformer ({\textsc BLT}), a novel large language model (LLM) architecture operating at the byte level, which, for the first time, achieves comparable performance to LLMs relying on traditional tokenization—while delivering notable gains in inference speed and robustness. {\textsc BLT} transforms byte sequences into variable-length patches that function as the model’s computational units. These patches are determined according to the predicted entropy of subsequent bytes, permitting greater computational attention and modeling resources in areas of heightened data complexity. We conduct the inaugural flop-controlled scaling analysis of byte-level architectures, extending to models with up to 8B parameters and exposed to 4T bytes during training. Our findings validate the practicality of building high-performing LLMs on raw bytes without resorting to predefined vocabularies. Both training and inference benefit from the adaptive use of longer patches for more predictable inputs, while empirical evaluations point to improved reasoning and better generalization on rare or unexpected cases. Collectively, for a fixed inference budget, {\textsc BLT} surpasses token-based models in scalability by allowing the simultaneous expansion of both patch length and model size.
\end{abstract}

\section{Introduction}
\label{section:intro}

We present the Byte Latent Transformer~(\textbf{BLT}), a novel tokenizer-free framework that processes raw byte sequences and, for the first time, achieves comparable performance to traditional tokenization-based systems at large scale, while delivering substantial gains in robustness and computational efficiency (\S\ref{sec:robustness}).
In contemporary large language models (llms), training is performed nearly end-to-end, but is still reliant on tokenization—a rule-based pre-processing stage that converts byte streams into a fixed vocabulary of tokens. This token selection influences string compression in a way that introduces several limitations, including dependency on the text's domain or modality~\citep{dagangetting}, increased vulnerability to input perturbations~(\S\ref{sec:robustness}), inadequate representation of orthographic features~\citep{edman2024cute}, and systemic disadvantages for some languages in multilingual applications~\citep{liang2023xlm,petrov2024language,limisiewicz2024myte}.

Tokenization has historically played a vital role because training llms directly on byte sequences incurs substantial computational expense, largely owing to increased sequence lengths~\citep{xue2022byt5}. Earlier research has sought to address this by leveraging more computationally efficient self-attention mechanisms~\citep{el2020characterbert, clark2022canine} or by exploring architectures that forgo attention entirely~\citep{wang2024mambabyte}~(\S\ref{section:related_work}). Nonetheless, these optimizations are chiefly beneficial for training \textit{small models}. When scaling up, the bulk of the computational burden in a Transformer arises from the extensive feed-forward network layers applied to each byte, rather than from the attention component itself.

In order to optimize compute usage, we introduce a dynamic, trainable approach for clustering bytes into \textit{patches}~(\S\ref{section:patching}) alongside a novel architecture that integrates both byte-level and patch-level features. In contrast to tokenization, {\textsc BLT} does not rely on a predetermined set of patch vocabulary. Instead, any combination of bytes can be projected into latent patch embeddings through compact, trainable encoder and decoder components. Our results demonstrate that this framework yields a \textit{greater} efficiency in compute distribution compared to models that depend on tokenization.

In conventional tokenization-based LLMs, each token is processed using an equal amount of computational resources. This design prioritizes overall performance but comes at the cost of efficiency, as tokens are generated via compression strategies that do not necessarily align with the varying complexity of different prediction tasks. The underlying principle in our proposed architecture is the adaptive allocation of computational effort depending on task difficulty. For instance, predicting the ending of most words typically constitutes a straightforward, low-entropy task, for which a large-scale transformer is superfluous; in contrast, selecting the initial token of a sentence requires greater modeling capacity. This conceptual framework informs the structure of {\textsc BLT}~(\S\ref{section:architecture}), which integrates three transformer blocks: two compact byte-level \textit{local models} alongside one expansive, global \textit{latent transformer}~(\autoref{fig:arch_high_level}).
To enable flexible grouping of bytes into patches and to optimize compute distribution, {\textsc BLT} segments input according to the entropy involved in predicting the subsequent byte. This results in context-sensitive byte clusters that maintain consistent information density.

We introduce the first study that systematically investigates how byte-level models scale under a fixed flop budget, training up to 8B parameter models and processing 4T bytes, without relying on fixed-vocabulary tokenization. Our results demonstrate that {\textsc BLT} achieves parity with Llama 3 in terms of training efficiency measured by flop consumption,\footnote{We assess a model’s computational expense by calculating the total number of Floating Point OPerations (flops) required.} yet requires as much as 50\% fewer flops for inference (\S\ref{section:scaling}).
Additionally, we provide evidence that leveraging raw bytes yields clear benefits in capturing infrequent patterns in the data. Compared to traditional tokenizer-based counterparts, {\textsc BLT} models show increased resilience to input noise and possess stronger capabilities at the character level, as verified on tasks involving orthographic reasoning, phonological phenomena, and low-resource machine translation~(\S\ref{sec:robustness}).
Moreover, by using {\textsc BLT} models, it becomes feasible to expand both model and patch sizes simultaneously while maintaining a constant inference flop allocation. Employing longer patch sizes typically reduces computation per forward pass, thereby freeing up capacity to expand the size of the global latent transformer, since it needs to be invoked less frequently. Our scaling studies—where inference flops are held constant (\autoref{fig:fixed_inference_scaling})—show that such byte-level models deliver substantially more favorable scaling behavior than traditional tokenization-dependent architectures.

In conclusion, the primary contributions of this work are as follows: 1) We propose {\textsc BLT}, a byte-latent large language model architecture capable of dynamically adjusting computational resources for greater flop efficiency, 2) Our experiments demonstrate that, up to the 8B parameter range, we can match Llama 3 in terms of training flops, with the flexibility to exchange small decreases in evaluation scores for substantial improvements in flop efficiency—reaching as much as 50\%, 3) {\textsc BLT} enables a new approach to llm scaling, allowing increases in model capacity while keeping inference computational costs constant, and 4) We provide evidence that {\textsc BLT} models are more resilient to noisy inputs and possess a heightened sensitivity to sub-word structure in data that token-based llms may overlook. The full training and inference code for {\textsc BLT} is publicly available at~\url{https://github.com/facebookresearch/blt}.

\section{Patching: From Individual Bytes to Groups of Bytes}
\label{section:patching}

Dividing byte sequences into \textit{patches} enables {\textsc BLT} to adaptively assign computational resources in response to input context. Figure~\ref{fig:patching_types} presents various techniques for partitioning bytes into patches. More precisely, a patching function $f_p$ partitions an input sequence of bytes $\pmb{x} = \{x_i | i = 1, \ldots, n\}$ of length $n$ into a sequence of $m < n$ patches $\pmb{p} = \{p_j | j = 1, \ldots, m\}$ by mapping each $x_i$ to the set \{0,1\}, where a 1 signifies the initiation of a new patch. Whether using token-based or patch-based architectures, the computational requirements are largely governed by the number of processing steps in the primary Transformer, which, for {\textsc BLT}, corresponds to the number of patches yielded by a specific patching function. Therefore, the principal determinant of processing expense—during both training and inference—is the mean size of patches, or \textit{patch size}, established by the chosen patching function~(\S\ref{section:flops}).
In what follows, we detail three specific patching strategies: using a uniform byte count per patch~(\S\ref{section:static-patch}), segmenting according to whitespace~(\S\ref{section:white-patch}), and adaptively patching guided by entropy estimates from a compact byte language model~(\S\ref{section:dyn-patch}). Additionally, we consider incremental patching methods and contrast patching with the concept of tokenization~(\S\ref{section:bpe-patch}).

\subsection{Strided Patching Every K Bytes}
\label{section:static-patch}

A common method for organizing bytes is to aggregate them into uniform patches of size $k$, as utilized in MegaByte~\citep{yu2023megabyte}. Employing a constant stride simplifies both the training and inference pipelines, allows for an intuitive means to modify the mean patch size, and offers direct control over the associated computational cost. Nevertheless, this approach has several major limitations. Primarily, it does not allocate computational resources adaptively based on local content—computational effort may be squandered on uninformative segments, such as whitespace in code (thereby underutilizing a transformer step $j$), or fall short in regions rich in information, like mathematical expressions. In addition, this rigid patching strategy fragments analogous byte sequences in a non-uniform and context-insensitive manner, potentially resulting in variable splits of identical words.

\subsection{Space Patching}
\label{section:white-patch}

\citet{slagle2024spacebyte} introduces an effective and straightforward enhancement to strided patching, whereby new patches are initiated immediately after the occurrence of any space-like byte\footnote{
    Space-like bytes refer to any byte that is neither a Latin character, a digit, nor a utf-8 continuation byte. Each patch is further required to include at least one byte that is not space-like. }. These boundaries often coincide with natural linguistic units in a range of languages. The method, termed Space patching, devotes a latent transformer computation (i.e., increased floating-point operations) to every word segment, guaranteeing a consistent patching approach to words across different sequences and ensuring computational effort is concentrated at points—such as immediately following spaces—where predictions are most challenging. For instance, given the question ``Who composed the Magic Flute?\uline{\hspace{.6em}}'', predicting the initial byte of the answer is significantly more complex than subsequent bytes following ``M'', as identifying the first character substantially narrows down the set of possible continuations, whereas completing ``Mozart'' after ``M'' is relatively straightforward. Nonetheless, Space patching is limited by its inability to handle all languages and domains effectively, and, crucially, by its inability to adapt patch sizes. In the following, we present a novel patching approach motivated by the observation that initial bytes in words are usually hardest to predict, but one that also allows flexible control over patch dimensions.

\subsection{Entropy Patching: Using Next-Byte Entropies from a Small Byte LM}
\label{section:dyn-patch}

Instead of utilizing a rule-based method like identifying whitespace, we adopt a data-centric strategy to detect locations where next-byte prediction uncertainty is high. Our method, termed \textit{entropy patching}, determines patch boundaries by leveraging entropy measurements.

We train a small byte-level auto-regressive language model on the training data for {\textsc BLT} and compute next byte entropies under the LM distribution $p_{e}$ over the byte vocabulary $\mathcal{V}$:

\begin{equation}
H(x_i) = \sum_{v \in \mathcal{V}} p_{e}(x_i=v|\pmb{x}_{<i}) \log p_{e}(x_i=v|\pmb{x}_{<i})
\end{equation}

We explore two approaches for detecting patch boundaries using entropies $H(x_i)$. The initial method selects points that exceed a global entropy threshold, as shown in~\autoref{fig:patching}. The alternative strategy looks for points with entropy values that are substantially higher than the preceding value. This method can also be viewed as spotting locations where the roughly monotonically decreasing entropy pattern within a patch is disrupted.

\begin{align*}
    \text{Global Constraint}\hspace{1cm}H(x_t)& > \theta_g\\
    \text{Approx. Monotonic Constraint}\hspace{1cm}H(x_t)& - H(x_{t-1}) > \theta_r
\end{align*}

Patch boundaries are determined through an efficient preprocessing process applied at the dataloading stage. This approach contrasts with the method in~\citet{nawrot-etal-2023-efficient}, where a classifier is employed to infer patch boundaries based on entropy measures. In the experiments we conduct~(\S\ref{section:experiments}), we evaluate and contrast these strategies for classifying bytes into low and high entropy categories.

\subsection{The Byte-Pair Encoding (BPE) Tokenizer and Incremental Patching}
\label{section:incremental-patching}
\label{section:bpe-patch}

Numerous contemporary llms, such as our baseline Llama 3, rely on subword tokenizers like bpe~\citep{gage1994new, sennrich-etal-2016-neural}. Throughout this work, ``tokens'' denote groups of bytes selected from a \textit{finite} pre-established vocabulary, whereas ``patches'' are used to describe adaptive groupings of bytes formed on-the-fly, independent of any preset vocabulary. One essential distinction between tokens and patches is that when utilizing tokens, the model does not have explicit access to the raw byte-level features.

An important advancement that {\textsc BLT} brings relative to tokenization-based architectures is a fundamental shift in how the balance between vocabulary size and computational resources is managed. In conventional llms, expanding the vocabulary results in tokens that represent more content on average, thereby reducing the total number of required inference steps; however, it simultaneously leads to a substantial increase in the output dimensionality of the final projection layer. This inherent compromise restricts tokenization-based frameworks from achieving considerable diversity in both token length and inference efficiency. As an illustrative example, Llama 3 enhances the mean token length from 3.7 to 4.4 bytes, but this improvement comes at the expense of quadrupling the size of its embedding matrix relative to Llama 2.

During generation, {\textsc BLT} must determine at each step of the byte sequence whether a patch boundary has been reached, since this affects whether the Latent Transformer is engaged for further computation. Crucially, this determination has to be made solely based on information available up to the current byte, as forthcoming parts of the sequence have not been generated yet. Therefore, patch segmentation cannot rely on future bytes. More precisely, we define a patching strategy $f_p$ to exhibit incremental patching if it holds that:
$$f_p(\pmb{x}_{<i}) = f_p(\pmb{x})_{<i}$$
The commonly used bpe algorithm does not possess this incremental property, because identical prefixes may be split into tokens differently depending on subsequent context; hence, it fails to meet the criterion above\footnote{Introducing a dedicated delimiter token to mark patch boundaries can adapt bpe to conform to incremental patching, but at the cost of lengthening the byte sequence.}.

\section{{\textsc BLT} Architecture}
\label{section:architecture}
The {\textsc BLT} framework consists of a sizable global autoregressive language model functioning on patch-level representations, as well as a pair of compact local models. These local models are responsible for converting byte sequences into patch representations and subsequently reconstructing the original bytes from those patch representations~(\autoref{fig:arch_high_level}).

\subsection{Latent Global Transformer Model}

\textit{The Latent Global Transformer} refers to an autoregressive transformer architecture $\mathcal{G}$ consisting of $l_{\mathcal{G}}$ layers, which processes a sequence of latent patch embeddings $p_j$ into corresponding output patch embeddings $o_j$. In this manuscript, we adopt the notation $j$ for patch indices and $i$ for byte indices. The global transformer is equipped with a block-causal attention mechanism~\citep{dubey2024llama} that limits each position’s attention scope to all preceding and current patches within a given document. Because this transformer accounts for the majority of computational operations during both pre-training and inference phases, selectively invoking the model enables us to dynamically allocate computational resources based on the intricacy of particular input and output sections.

\subsection{Local Encoder}
\label{section:local}

The \textit{Local Encoder Model}, represented as $\mathcal{E}$, employs a compact transformer structure characterized by a significantly reduced number of layers, $l_{\mathcal{E}} << l_{\mathcal{G}}$, focusing on the efficient conversion of input byte sequences $b_i$ into rich patch representations $p_j$. Notably diverging from traditional transformer models, it incorporates a cross-attention layer subsequent to each transformer layer to aggregate byte-level information into higher-level patch representations~(\autoref{fig:crossattn}).

Initially, each byte $b_i$ within the sequence undergoes embedding via a $\mathbb{R}^{256 \times h_{\mathcal{E}}}$ matrix, yielding vectors $x_i$. These base embeddings may be further enhanced through the integration of hash-embeddings~(\S\ref{sec:hash-emb}). 

The model then alternates between transformer layers and cross-attention mechanisms~(\S\ref{sec:enc-cross-attn}), progressively evolving these initial byte-level representations into patch-level representations $p_i$, which are subsequently delivered to the global transformer component, $\mathcal{G}$, for further processing.

Within this architecture, each transformer's attention mechanism is configured with a \textit{local block causal} attention mask, restricting the receptive field for any given byte to the $w_{\mathcal{E}}$ previous bytes. This windowed attention may span over dynamically defined patch boundaries, but does not extend across distinct document boundaries. The subsequent sections elaborate on the construction of the embedding vectors and the design of the cross-attention block.

\subsubsection{Encoder Hash n-gram Embeddings}
\label{sec:ngram-emb}
\label{sec:hash-emb}
An essential aspect of developing strong and informative representations at each position $i$ involves encoding contextual information from preceding bytes. Within {\textsc BLT}, this is accomplished by representing the current byte $b_i$ on its own as well as within the context of a byte n-gram. Specifically, for each step $i$, we begin by generating byte-grams

\begin{align}
    g_{i,n}=\{b_{i-n + 1},\ldots, b_i\}
\end{align}

for every byte position $i$ and for each $n$ in the range from three to eight.\footnote{
Byte-grams whose size $n$ exceeds the current index $i$ (i.e., $i < n$) are not considered.
}

Next, we propose hash $n$-gram embeddings, where each byte $n$-gram is assigned an index within an embedding table $E_{n}^{hash}$ by applying a hash function, with $E_{n}^{hash}$ maintaining a fixed size for each $n \in \{3,4,5,6,7,8\}$, following~\citep{bai2010learning}. 
The embedding obtained for the byte $n$-gram is then incorporated into the corresponding byte embedding. Afterwards, normalization is performed before inputting the representation to the local encoder. The computation for the enhanced embedding is as follows:

\begin{align}
e_i &= x_i + \sum_{n = 3,...,8}E_{n}^{hash}(\text{Hash}(g_{i,n})) \\
\text{where, Hash}(g_{i,n}) &= \text{RollPolyHash}(g_{i,n}) \% |E_{n}^{hash}|
\end{align}

Each $e_i$ is scaled by dividing by the total number of $n$-gram sizes plus one, and we apply RollPolyHash as specified in Appendix~\ref{appendix:rollpolyhash}. Section~\ref{section:ablations} investigates the influence of $n$-gram hash embeddings by varying $n$ and the size of the embedding table, analyzing their impact on scaling law behavior under a fixed computational budget. Beyond the use of hashed $n$-gram embeddings, we also explored frequency-based $n$-gram embeddings; further information on these experiments is given in Appendix~\ref{appendix:ngrams}.

\subsubsection{Encoder Multi-Headed Cross-Attention}
\label{sec:enc-cross-attn}
Our design largely adopts the input cross-attention approach from the Perceiver architecture~\citep{jaegle2021perceiver}, but differs in a key aspect: instead of mapping to a static collection of latent vectors, the latent representations in our model are linked to the specific patch representations, as shown in~(\autoref{fig:crossattn}). Each latent vector therefore focuses exclusively on the bytes contained within its respective patch. The process begins with forming a query vector for every patch $p_j$ by aggregating the byte embeddings of patch $p_j$ (typically via pooling), then subjecting this pooled vector to a linear transformation, $\mathcal{E}_{C} \in \mathbb{R}^{h_{\mathcal{E}} \times (h_{\mathcal{E}} \times U_{\mathcal{E}})}$, where $U_{\mathcal{E}}$ stands for the number of attention heads in the encoder cross-attention mechanism. Explicitly, letting $f_{\text{bytes}}(p_j)$ be the ordered bytes of patch $p_j$, we then compute

\begin{align}
P_{0,j} &= \mathcal{E}_{C}(f_\text{bytes}((p_j)), f~ \text{is a pooling function} \\
P_l &= P_{l-1} + W_o\left(\text{softmax}\left(\frac{QK^T}{\sqrt{d_k}}\right)V\right) \\
\text{where } Q_j &= W_q(P_{l-1,j}), K_i = W_k(h_{l-1,i}), V_i = W_v(h_{l-1,i}) \\
h_l &= \text{Encoder-Transformer-Layer}_l(h_{l-1})
\end{align}

Here, $P \in \mathbb{R}^{n_p \times h_{\mathcal{G}}}$ denotes the set of $n_p$ patch representations intended for processing by the global model. These are initially formed by aggregating the byte-level embeddings $e_i$ that make up each patch $p_j$. The matrices $W_q$, $W_k$, $W_v$, and $W_o$ are linear projection weights for the queries, keys, values, and output, respectively; the keys and values are generated by projecting the byte representations $h_i$ from the preceding layer (or by using $e_i$ in the initial layer). For patch-based masking, we adopt a strategy wherein each query $Q_j$ is restricted to attending solely to those keys and values that are associated with bytes from the same patch $j$. As the architecture leverages multi-head attention over $Q, K,$ and $V$, and since patch representations generally possess a higher dimensionality ($h_{\mathcal{G}}$) compared to $h_{\mathcal{E}}$, we structure $P_l$ as a set of heads each of size $h_{\mathcal{E}}$ during the cross-attention operation and subsequently concatenate these head outputs to form representations of size $h_{\mathcal{G}}$. Furthermore, the design incorporates pre-LayerNorm normalization for queries, keys, and values, and omits the use of positional embeddings in the cross-attention mechanism. The output of this cross-attention block is finally combined with its input via a residual connection.

\subsection{Local Decoder} 

In parallel with the design of the local encoder, the local decoder $\mathcal{D}$ also utilizes a compact transformer-based architecture, comprising $l_{\mathcal{D}} << l_{\mathcal{G}}$ layers. Its purpose is to map the input sequence of global patch embeddings $o_j$ back to a stream of raw bytes, denoted as $y_i$. The process of byte prediction by the local decoder is inherently sequential, with each prediction conditioned on the preceding decoded bytes. Accordingly, it accepts as input the hidden states generated by the local encoder for the byte sequence. Structurally, it alternates through $l_{\mathcal{D}}$ layers, where each block consists of a cross-attention mechanism followed by a transformer layer. In this configuration, cross-attention is employed prior to the transformer component to initially translate patch representations into corresponding byte-level representations, after which the transformer's attention layers refine the resulting byte sequence.

\subsubsection{Decoder Multi-headed Cross-Attention} 

In the decoder cross-attention, the roles of the queries and key/values are interchanged i.e. the byte-representations are now the queries, and the patch representations are now the key/values. The initial byte-representations for the cross-attention are initialized as the byte embeddings from the last encoder layer i.e. $h_{l_{\mathcal{E}}}$. The subsequent byte-representations for layer $l$, $d_{l,i}$ are computed as:

\begin{align}
D_0 &= h_{l_{\mathcal{E}}} \\
B_l &= D_{l-1} + W_o\left(\text{softmax}\left(\frac{QK^T}{\sqrt{d_k}}\right)V\right), \\
  &\text{where } Q_i = W_q(d_{l-1,i}), K_i = W_k(\mathcal{D}_C(o_j)), V_i = W_v(\mathcal{D}_C(o_j)) \\
  D_l &= \text{Decoder-Transformer-layer}_l(B_l)
\end{align}

Here, $W_k$ and $W_v$ denote the projection matrices used to obtain keys and values, respectively, which are applied to the linearly transformed and split outputs $\mathcal{D}_C$—derived from the global model’s final patch representations $o_j$. The matrix $W_q$ serves as the projection for queries, operating on the byte-level representations $d_{l-1}$ from the previous decoder transformer block (or on $h_{l_{\mathcal{E}}}$ in the case of the first decoder layer), while $W_o$ is responsible for projecting outputs. As a result, $B$ has the shape $B \in \mathbb{R}^{h_{\mathcal{D}} \times n_b}$, where $n_b$ denotes the total output byte count. The next-step decoder representations, $D_l$, are generated by applying a decoder transformer layer to the cross-attention output $B$. Consistent with the approach used in the local encoder cross-attention, the mechanism employs multiple attention heads, incorporates pre-LayerNorm, excludes positional embeddings, and includes a residual connection surrounding the cross-attention module.

\section{Experimental Setup}
\label{section:experiments}
We construct rigorously controlled experiments to assess {\textsc BLT} against models employing tokenization, making explicit efforts to ensure that {\textsc BLT} does not benefit from access to extended sequence contexts.

\subsection{Pre-training Datasets} For all model scales analyzed in this study, pre-training was conducted using two main datasets: 1) The Llama 2 dataset~\citep{touvron2023llama}, consisting of 2 trillion tokens amassed from diverse publicly available sources, with extensive cleaning and filtering applied to enhance dataset quality; and 2) {\textsc BLT}-1T: a newly curated dataset containing 1 trillion tokens, sourced from an array of public domains and supplemented with selected portions of the Datacomp-LM pre-training set~\citep{li2024datacomp}. The Llama 2 dataset is leveraged in scaling law investigations, specifically regarding the optimal token count as identified by~\cite{dubey2024llama}, to inform the architectural choices for {\textsc BLT}. The {\textsc BLT}-1T dataset underpins the full pre-training process, enabling comprehensive downstream comparisons with Llama 3. Importantly, both datasets exclude any data derived from Meta platforms or offerings. Additionally, for tokenizer-based baseline assessments, the Llama 3 tokenizer—characterized by a 128K token vocabulary—was adopted, as it yielded better baseline results compared to the Llama 2 tokenizer across our experiments.

\subsection{Entropy Model} For our experiments, the entropy model is implemented as a byte-level language model, utilizing the same training data distribution as the complete {\textsc BLT} system. Unless specified otherwise, the architecture selected is a transformer comprising 100M parameters, 14 layers, with a hidden dimension of 512, and a sliding window attention span of 512 bytes. All additional hyperparameters are matched to those in our local and global transformer configurations. A range of model sizes, receptive fields, and architectural choices were evaluated, as described in \autoref{section:ablations}. Notably, with sufficiently limited receptive fields, the resulting entropy model allows for an efficient encoding as a lookup table.

\subsection{Entropy Threshold and Equalizing Context Length} In entropy-based patching approaches, we determine an appropriate entropy threshold such that the resulting average \textit{patch size} on the pretraining dataset aligns with our target. Within {\textsc BLT}, as opposed to conventional tokenization, the \textit{patch size} is freely adjustable, which substantially affects the amount of contextual information accessible to the model. To prevent models with larger patch sizes from being inadvertently favored due to increased context, we fix the expected total number of bytes per batch across all experimental setups. Consequently, the input sequence length is decreased for models with greater patch sizes. Specifically, for the Llama 2 dataset, we set the context window to 8k bytes, whereas for the {\textsc BLT}-1T corpus, we expand the context to an average of 16k bytes, holding the average batch size constant at 16M bytes in both cases.

Although the mean batch size remains unchanged, utilizing dynamic patching approaches results in varying proportions of bytes per patch during data batching. To improve efficiency, our {\textsc BLT} training setup compacts batches by grouping patches tightly, thereby eliminating the need for padding in the computationally intensive latent transformer. This strategy maintains a consistent count of patches across all batches. To control memory usage and prevent surges caused by sequences containing exceptionally large patches, we pad or, if necessary, truncate input byte sequences to 12k bytes for Llama 2 and 24k bytes for {\textsc BLT}-1T datasets during the training process.

\subsection{Entropy Model Context}
\label{section:entropy-context}
Our experiments demonstrate that applying entropy patching results in increasingly larger patches, particularly in tasks featuring structured formats such as multiple choice (refer to patching results on an MMLU item in \autoref{fig:mmlu_patching}), where the content exhibits high repetition. These observed differences are attributable to reduced entropy levels within the repeated elements present in the entropy model context. Therefore, for the extensive evaluation using {\textsc BLT}-Entropy at a patch size of 4.5, we reinitialize the entropy context through line breaks and incorporate an approximate monotonicity constraint, since this approach mitigates issues arising from "entropy drift" due to changing context length. While this modification alters the method by which entropies are estimated, our procedure for determining the entropy threshold value remains unchanged.

\subsection{FLOPs Estimation} 
\label{section:flops}

Our approach closely adheres to the methodology outlined in Chinchilla~\citep{hoffmann2022training} for estimating transformer floating point operations (flops), accounting for contributions from the feed-forward layers, qkvo projections within self-attention, and both attention calculation and output projection. However, we diverge from this framework by considering the input embedding layer as an efficient lookup rather than a dense matrix multiplication, effectively reducing its computational cost to zero flops. Consistent with established literature, we assume that the backward pass requires double the number of flops relative to the forward pass.

To compute flops \textit{per byte} for {\textsc BLT} models, we add up the flops for the local encoder transformer, the global latent transformer, and the local decoder transformer, together with the cross attention blocks in the encoder and the decoder:

\begin{align}
    \text{FL}_{\text{{\textsc BLT}}} &= \text{Transf. FL}(h_{\mathcal{G}}, l_{\mathcal{G}}, m=n_{ctx}/n_p, V=0) / n_p\\
   &+ \text{Transf. FL}(h_{\mathcal{E}}, l_{\mathcal{E}}, m=w_{\mathcal{E}}, V=0) \\
   &+ \text{Transf. FL}(h_{\mathcal{D}}, l_{\mathcal{D}}, m=w_{\mathcal{D}}, V=256) \\ 
    &+ \text{Cross Attn. FL}(h_{\mathcal{E}}, l_{\mathcal{E}}, m=n_p, r=n_p/k) \times k/n_p \\ 
   &+ \text{Cross Attn. FL}(h_{\mathcal{D}}, l_{\mathcal{D}}, m=k, r=k/n_p) 
\end{align}

Here, $n_{ctx}$ represents the length of the sequence in bytes, while $n_p$ denotes the patch size. The parameter $r$ specifies the proportion of queries relative to key/values, and $k$ defines the patch-to-byte dimensional ratio, that is, the count of local model partitions concatenated to generate the overall model embedding ($k=2$ in \autoref{fig:crossattn}). $V$ stands for the output projection vocabulary size, utilized solely within the local decoder. The sequence—whether processed at the byte-level or patch-level—influences the attention’s context length, denoted as $m$. We revise the attention flop calculations based on the context for each module. The detailed flop computation formulae for both Transformer-FLOPs and Cross-Attention FLOPs can be found in Appendix~\ref{section:flops-appendix}.

\subsection{Bits-Per-Byte Estimation} Perplexity only makes sense in the context of a fixed tokenizer as it is a measure of the uncertainty for each token. When comparing byte and token-level models, following previous work~\citep{xue2022byt5, yu2023megabyte, wang2024mambabyte}, we instead report Bits-Per-Byte (BPB), a tokenizer independent version of perplexity. Specifically:

\begin{align}
    \text{BPB}(x)&= \frac{\mathcal{L}_{CE}(\pmb{x})}{\ln(2)\cdot n_{\text{bytes}}}
\end{align}

Here, the cumulative cross-entropy loss quantifying the uncertainty for data $\pmb{x}$ is divided by both the total byte count of $\pmb{x}$ and a constant, providing a normalized measure.

\subsection{Transformer Architecture Hyperparameters} The transformer blocks used throughout {\textsc BLT}, encompassing both the local and global components, are based primarily on the Llama~3~\citep{dubey2024llama} framework. Within the feed-forward networks, the SwiGLU activation function~\citep{swiglu} is employed. Self-attention modules incorporate rotary positional embeddings (RoPE)~\citep{RoPe} configured with $\theta=500000$ as specified in~\citep{xiong2024effective}. For normalization across layers, we utilize RMSNorm~\citep{RMSNorm}. Regarding self-attention layers utilizing conventional attention masks—such as those for \textit{block causal} or \textit{fixed-window block causal} attention—Flash attention~\citep{dao2022flashattention} is adopted, implementing a fixed window length of 512 for windowed attention. As the cross-attention modules require dynamic, patch-based masking, Flex Attention\footnote{\url{https://pytorch.org/blog/flexattention}} is leveraged to yield fused operator implementations and provide considerable acceleration during training.

\subsection{{\textsc BLT}-Specific Hyperparameters} 
To evaluate the performance of {\textsc BLT} models, we perform experiments in two primary areas: examining scaling behaviors and assessing results on downstream tasks. These experiments span model sizes of 400M, 1B, 2B, 4B, and 8B parameters. Details regarding architecture-specific hyperparameters for each configuration can be found in Appendix~Table~\ref{tab:arch}.

For the local encoder, we initialize the queries for the initial cross-attention layer using max-pooling. All {\textsc BLT} variants use $500,000$ hashes and a single hash function; n-gram sizes are selected from 3 to 8. The learning rate is consistently set to $4e-4$ across all experiments. This learning rate was identified as optimal based on a sweep over $1e-3$ to $1e-4$ on both 400M and 1B model sizes, where the performance of token and {\textsc BLT} models indicated no preference for different values. For scaling analyses leveraging the Llama-2 dataset, we adopt batch-sizes as advised in~\cite{dubey2024llama}, or their corresponding values in bytes when applicable. The optimization process employs the AdamW optimizer~\citep{adamw} with $\beta_1$ fixed at 0.9, $\beta_2$ at 0.95, and $\epsilon$ of $10^{-8}$. We utilize a linear warm-up over 2000 steps and apply a cosine annealing schedule for learning rate decay to zero. Weight decay is set at 0.1, and gradient norms are globally clipped at a maximum of 1.0.

\section{Scaling Trends}
\label{section:scaling}

We provide a comprehensive analysis of scaling behaviors for byte-level models, offering insights relevant to the further development of {\textsc BLT} models. Our investigation into scaling seeks to overcome constraints found in earlier studies on byte-level models by: (a) examining scaling patterns within the context of compute-optimal training approaches, (b) training 8B parameter models with substantial data volumes—reaching up to 1T tokens or 4T bytes—and conducting evaluations across a range of downstream applications, and (c) assessing scaling patterns under conditions where inference cost is held constant. Subsequent sections will delve into the unique benefits associated with modeling sequences at the byte level.

\subsection{Parameter Matched Compute Optimal Scaling Trends}
\label{section:parameter-matched-scaling}
We utilize the Llama 2 dataset to train a series of \textit{compute-optimal} bpe and {\textsc BLT} models at four distinct model sizes, ranging from 1B to 8B parameters. Subsequently, we graph the amount of training flops required versus the language modeling performance, evaluated on a representative sample of the training data mix. For bpe models, training occurs with the optimal proportion of model parameters relative to the amount of training data, in accordance with the findings in Llama~3~\citep{dubey2024llama}. This \textit{compute-optimal} configuration is, in theory, engineered to maximize performance on the training data for a fixed computational budget~\citep{hoffmann2022training}, thereby establishing a solid benchmark for comparison. Corresponding {\textsc BLT} models are also trained for each bpe model, employing a Latent Transformer that is architecturally and parametrically identical to its bpe Transformer counterpart, and utilizing the same dataset.

As shown in~\autoref{fig:scaling-compute-opt} (right), {\textsc BLT} models consistently achieve performance comparable to or better than their bpe-based counterparts, with this pattern persisting across increases in model size and computational resources (flops). According to our knowledge, {\textsc BLT} represents the first byte-level Transformer to demonstrate comparable scaling behavior to BPE models in compute-efficient training regimes. These observations support our hypothesis that the optimal balance between parameter count and training computation identified for bpe models is also applicable, or at least closely similar, for {\textsc BLT}.

Achieving parity with bpe scaling patterns requires both enhancements in model architecture and the use of dynamic patching strategies. As depicted in ~\autoref{fig:scaling-compute-opt} (left), we contrast models that utilize space-patching with Llama 3, where our approximation of SpaceByte \citep{slagle2024spacebyte} relies on {\textsc BLT} space-patching, omitting n-gram embeddings and cross-attention modules. While SpaceByte represents an improvement relative to Megabyte, it still lags behind Llama 3 by a considerable margin. In \autoref{fig:scaling-compute-opt} (right), we demonstrate the impact brought by integrating both updated architectural features and dynamic patching mechanisms. At scale, {\textsc BLT} models match the performance of leading tokenizer-based models such as Llama 3.

Additionally, we find that for models utilizing a tokenizer, the selection of tokenizer significantly influences model performance; specifically, models trained with the Llama-3 tokenizer achieve better results than those trained with the Llama-2 tokenizer when both are exposed to identical training data.

In summary, our {\textsc BLT} framework exhibits scaling patterns that are intermediate between Llama 2 and Llama 3 when considerably increasing patch size. Whereas the bpe tokenizers for Llama 2 and 3 have mean token lengths of 3.7 and 4.4 bytes respectively, {\textsc BLT} attains analogous scaling properties with mean patch sizes of 6 and even 8 bytes. Because inference flops are inversely related to average patch size, adopting an 8-byte patch size can result in approximately a 50\% reduction in inference flops. Additionally, models employing larger patch sizes appear to benefit more as both model and dataset scales increase. Notably, while {\textsc BLT} with an 8-byte patch size initially underperforms relative to bpe Llama 2 at the 1B scale, it surpasses bpe at the 7B scale. This finding implies that such large patch sizes could deliver enhanced performance at greater scales and raises the possibility that even larger patch sizes might be advantageous as both model capacity and training resources expand.

\subsection{Beyond Compute Optimal Task Evaluations}
\label{section:evals}
To further analyze scaling behavior, we extend the training of an 8B {\textsc BLT} model past the compute-optimal regime, employing the {\textsc BLT}-1T dataset, which is both larger and of higher quality. The resulting model is then evaluated on a range of widely-used classification and generation tasks. For these assessments, we focus on tasks targeting common sense reasoning, world knowledge, and code generation:

\paragraph{Classification tasks} encompass ARC-Easy (0-shot)~\citep{Eval_ARC}, ARC-Challenge (0-shot)~\citep{Eval_ARC}, HellaSwag (0-shot)~\citep{Eval_hellaswag}, PIQA (0-shot)~\citep{Eval_piqa}, and MMLU (5-shot)~\citep{hendrycks2020measuring}. The evaluation utilizes prompt-scoring, where we compute the probability of choice-specific tokens and subsequently derive mean accuracy across these tasks.

\paragraph{Coding related generation tasks:} We measure pass@1 rates on MBPP (3-shot)~\citep{Eval_MBPP} and HumanEval (0-shot)~\citep{Eval_HumanEval} datasets to assess the capability of LLMs in generating Python code.

In \autoref{tab:evals}, we present a comparison among three different models that were trained on the {\textsc BLT}-1T dataset: a model leveraging the Llama 3 tokenizer with byte pair encoding (bpe),\footnote{We opt for the Llama 3 tokenizer and its 128k vocabulary due to its superior performance relative to the 32k vocabulary of Llama 2.} as well as two alternative versions of the {\textsc BLT} architecture. These variants include {\textsc BLT}-Space, which adopts a space-patching method, and {\textsc BLT}-Entropy, which implements an entropy-guided patching approach. In the case of the entropy-based method, we impose an approximate monotonicity constraint and refresh the model’s context at line breaks, following the approach described in~\autoref{section:entropy-context}. Training for all three models was conducted using a matching FLOP allocation to ensure a fair comparison. For {\textsc BLT}-Entropy, we further adjust the entropy threshold at inference time—lowering it from 0.6 to 0.1—discovering that this change enhances task outcomes, although it also results in an increased number of inference steps.

The {\textsc BLT}-Entropy model surpasses the Llama 3 model in performance on 4 of the 7 evaluated tasks, even though both models are trained using an equivalent amount of data. This advantage can likely be attributed to (1) the more efficient utilization of computational resources through dynamic patching, and (2) the explicit modeling of information at the byte level rather than the token level.

Conversely, {\textsc BLT}-Space generally lags behind the Llama 3 tokenizer in all but one evaluated task. Nonetheless, its use of a larger average patch size of 6 bytes brings a notable decrease in inference flops. For context, both bpe and entropy-patching models display similar average patch sizes, around 4.5 bytes, across the training dataset mixture. When subjected to an identical training budget, the model utilizing the larger patch size processes approximately 30\% more data than the bpe and entropy-patching counterparts, which may result in {\textsc BLT} deviating further from the compute-optimal threshold.

\subsection{Patches Scale Better Than Tokens} With {\textsc BLT} models, it is feasible to expand both the model and patch dimensions in tandem, all while preserving a consistent training and inference computational budget and utilizing the same amount of training data. Patch-based architectures uniquely enable increasing patch size beyond traditional limitations, sidestepping the efficiency constraints encountered in fixed-vocabulary token-based models, as elaborated in Section~\ref{section:bpe-patch}. Employing longer patches leads to computational savings that can be redirected to augment the capacity of the global latent transformer, since the transformer is invoked less frequently.

We perform a fixed inference scaling analysis to evaluate whether increasing model size while reducing the number of steps on larger patches yields superior performance compared to employing smaller models that process more steps. Utilizing initial model configurations of 400 million and 3.6 billion parameters with the Llama 2 tokenizer, we identify flop-matched models equipped with the Llama 3 tokenizer, as well as {\textsc BLT}-Entropy architectures employing mean patch sizes of 6 and 8 bytes on the training datamix (refer to~\autoref{tab:fixed-inf-params} for further specifications). For models operating with a patch size of 8, we implement 3 encoder layers rather than 1. Each model is subsequently trained under several distinct training flop budgets.

As depicted in \autoref{fig:fixed_inference_scaling}, {\textsc BLT} architectures exhibit superior scaling behavior compared to tokenization-based models across both categories of inference computation. While BPE-based models initially outperform at lower training resource levels, their performance is soon overtaken by that of {\textsc BLT} as compute increases slightly beyond the regime optimal for their size. Consequently, allocating greater resources for a one-off pretraining phase can be advantageous, resulting in more capable models at a constant inference computational cost. This pattern is clearly exemplified by the 8B model class, such as Llama 3.1, which underwent training on datasets approximately two orders of magnitude larger than the compute-optimal scale corresponding to its size.

The point at which {\textsc BLT} surpasses token-based models in performance has moved somewhat nearer to the compute-optimal threshold with the adoption of higher flop class models (from 3x to 2.5x the optimal compute allocation). Likewise, the model with a larger patch size of 8 exhibits a more rapid scaling behavior in these larger flop settings, overtaking other variants earlier in the process. As indicated in Section~\ref{section:parameter-matched-scaling}, models with larger patch sizes tend to approximate the performance of BPE models more closely at increased model scales. This phenomenon can be partly explained by the relatively reduced proportion of total computational resources consumed by the byte-level Encoder and Decoder components, which do not scale as rapidly as the Latent Transformer. For instance, scaling the total parameter count by a factor of 20—from 400M to 8B—leads to only an approximate twofold increase in local parameters for {\textsc BLT}. This distinction is notable, as increased patch sizes influence the computational cost of only the patch Latent Transformer, leaving the byte-level modules largely unaffected. Indeed, this explains the observation that the model size ratio for {\textsc BLT}-Entropy with patch size 8 rose from 1.6x to 1.7x that of the Llama 2 baseline as the model scale increased.

In conclusion, our investigation into scaling with patch length reveals that the {\textsc BLT} patch-based architecture exhibits improved scaling behavior when both patch size and overall model size are increased together. These positive scaling effects continue and, in fact, become more pronounced as model sizes grow further.

\section{Byte Modeling Improves Robustness} We further evaluate the robustness of {\textsc BLT} relative to token-based counterparts that do not incorporate explicit byte-level input, and introduce a technique for transforming pretrained token-based models to operate at the byte level.
\label{sec:robustness}

\subsection{Character-Level Tasks}

One of the initial incentives behind developing byte-level models stemmed from their inherent resilience to noise at the byte level within inputs, in addition to their capacity to recognize token structure—something that traditional tokenizer-based approaches typically lack. To assess these properties, we conduct further experiments using benchmark datasets specifically designed to gauge robustness against noisy inputs and to test sensitivity to characters in both English and various other languages, covering digits and phonemes as well. The findings of these experiments are summarized in Table~\ref{tab:char_tasks}.

\paragraph{Noisy Data} To evaluate the robustness of {\textsc BLT} relative to tokenizer-based models, we generate modified versions of the benchmark classification datasets introduced in Section~\ref{section:evals}. Five distinct character-level perturbation techniques are used to create these noisy datasets: (a) \textit{AntSpeak}: Reformats text by rendering it entirely in uppercase letters with spaces inserted between each character. (b) \textit{Drop}: Randomly eliminates 10\% of the characters within the text sequence. (c) \textit{RandomCase}: Arbitrarily changes half of the characters to uppercase and the remaining half to lowercase. (d) \textit{Repeat}: Selects 20\% of the characters at random and repeats each of them up to four times. (e) \textit{UpperCase}: Converts all characters in the text string to uppercase. When conducting the evaluation, each noising procedure is independently applied to the prompt, the completion, or both, treating each as a distinct condition, and the mean score across these variants is computed. Table \ref{tab:char_tasks} displays the results on a noised version of HellaSwag~\citep{Eval_hellaswag}, demonstrating that {\textsc BLT} consistently surpasses tokenizer-dependent models in terms of robustness, achieving an average lead of 8 points over the model trained on an identical dataset and even outperforming the Llama 3.1 model, which was trained on a substantially larger corpus.

\paragraph{Phonology - Grapheme-to-Phoneme (G2P)} We evaluate the extent to which {\textsc BLT} can convert a sequence of graphemes (letters denoting a word) into the phoneme sequence representing its pronunciation. Table \ref{tab:char_tasks} summarizes our findings on the G2P task under a 5-shot regime using Phonology Bench~\citep{suvarna2024phonologybench}, showing that {\textsc BLT} achieves superior performance relative to the baseline Llama 3 1T model that relies on tokenizer-based processing.

\paragraph{CUTE}
To evaluate the character-level capabilities of {\textsc BLT}, we employ the CUTE benchmark~\citep{edman2024cute}, which encompasses a diverse range of tasks grouped into three principal categories: compositional understanding, orthographic similarity, and sequence manipulation. This benchmark is particularly demanding for most models utilizing tokenization, as although such models generally have implicit access to token spellings, they typically struggle to leverage this information for practical text manipulation. As detailed in Table~\ref{tab:char_tasks}, {\textsc BLT}-Entropy achieves a performance exceeding both BPE Llama 3 models by more than 25 points on these tasks. Notably, our model excels in tasks involving character-level manipulations, registering 99.9\% accuracy on spelling-related subtasks. These marked gains are achieved even though {\textsc BLT} is trained on 16 times less data compared to Llama 3.1, suggesting that BPE architectures encounter substantial difficulty in acquiring robust character-level representations. Figure \ref{fig:cute_examples} presents selected cases where the Llama 3 tokenizer-based model fails, whereas our {\textsc BLT} model succeeds. Tasks requiring word deletion and insertion remain the only areas where BPE models outperform; such operations may pose greater complexity for a byte-level system, although the disparity is limited and building up from characters to words may ultimately prove more tractable than breaking down from words to characters. Across all experiments, we employ a consistent evaluation protocol and utilize the original Huggingface prompts. Enhanced prompt engineering could potentially improve outcomes for BPE models.

\paragraph{Low Resource Machine Translation} We assess the performance of {\textsc BLT} on translation tasks involving six major language families and twenty-one lower resource languages, each utilizing diverse scripts, as provided by the FLORES-101 benchmark~\citep{goyal2022flores}. SentencePiece BLEU scores for these tasks are displayed in Table~\ref{tab:flores}. The findings indicate that {\textsc BLT} delivers superior results compared to a model utilizing the Llama 3 tokenizer, exhibiting an overall gain of 2 BLEU points when translating into English and a 0.5-point increase when translating out of English. For widely spoken language pairs, {\textsc BLT} matches or marginally surpasses Llama 3 in performance. Crucially, {\textsc BLT} demonstrates marked improvements over Llama 3 on many language pairs representing lower-resource families, highlighting the strength of byte modeling in handling diverse and uncommon byte sequences.

\subsection{Training {\textsc BLT} from Llama 3}
\label{sec:distillation}
We investigate an approach in which {\textsc BLT} models benefit from previously trained tokenizer-based models, enhancing both the efficiency and effectiveness of convergence during training. Specifically, we initialize the global transformer parameters of {\textsc BLT} using those from a pre-trained Llama 3.1 model. Following this initialization, we proceed to fine-tune the global transformer parameters with a learning rate that is one-tenth of that applied to the local encoder and local decoder, and maintain the number of optimization steps typically recommended for Llama 3. Table~\ref{tab:distill} provides a comparative analysis against a baseline {\textsc BLT}. The results reveal that {\textsc BLT} initialized from Llama 3.1 outperforms both standard Llama 3 and baseline {\textsc BLT}—despite equivalent computational budgets. Additionally, relative to our {\textsc BLT}-Entropy variant (shown in Table~\ref{tab:evals}), which is trained on a much larger corpus of 1T tokens, {\textsc BLT} initialized from Llama 3.1 still surpasses it on the MMLU benchmark. This finding highlights the promise of this initialization strategy for substantially lowering the necessary training FLOPs.

This approach can be interpreted as adapting tokenizer-dependent architectures to operate without tokenizers, thus reconfiguring a pre-trained LLaMA 3.1 model into a {\textsc BLT} model. For thorough comparison, we present results for both the original LLaMA 3.1—trained on 15T tokens—and the {\textsc BLT} model derived from LLaMA 3 in Table \ref{tab:distill}. While the adapted model maintains nearly equivalent outcomes on MMLU and HumanEval, it exhibits a more pronounced decrease in performance across various other benchmarks. These observations indicate that further optimization—particularly with regard to selecting appropriate data compositions and tuning hyperparameters—is required to more effectively utilize the pre-trained model and enhance overall results.

\section{Ablations and Discussion}
\label{section:ablations}
Here, we present ablation studies to support the architectural decisions made for {\textsc BLT}, as well as the choices regarding the patching methodology and hyper-parameter selection employed in training the {\textsc BLT} model with 8B parameters on the {\textsc BLT}-1T dataset.

\paragraph{Entropy Model Hyper-parameters} In order to analyze how entropy model capacity and the length of the context window influence scaling behavior, we conduct experiments training byte-level entropy transformer architectures across a spectrum of model sizes ranging from 1 million to 100 million parameters and evaluate a range of context window lengths between 64 and 512. We generate scaling law plots of bits per byte (bpb) against training flop count, using our $400m$ and $1b$ {\textsc BLT} models trained on the Llama-2 corpus, as illustrated in \autoref{fig:entropy_ablation}. The results indicate that larger model size and longer context windows both improve scaling outcomes, although performance gains become marginal beyond 50 million parameters.

\paragraph{Types of Patching}

We conduct an ablation study of the four patching strategies described in Section~\ref{section:patching}: (1) Strided Patching employing strides of 4 and 6, (2) Whitespace-based patching, (3) Patching utilizing the BPE Tokenizer from the Llama 3 model, and (4) Entropy-guided patching leveraging a compact byte-level language model.

Although dynamic patching shortens the actual sequence length, we ensure that the context window remains comparable across all patching strategies by standardizing the sequence length. Consequently, each model processes an equivalent number of bytes per sequence during both training and inference on average, minimizing potential confounds related to modeling longer contexts. The outcomes of these comparative analyses are presented in Figure~\ref{fig:scaling-compute-opt}. All alternative patching methods yield better performance than static patching, with space patching exhibiting results nearly matching those of dynamic entropy-based patching.

As shown in \autoref{tab:evals_ablation}, we evaluate the performance of {\textsc BLT} models using different tokenization and patching methods, namely tokenizer-based approaches, space patching, and entropy-based patching. All models are trained on the Llama 2 dataset for an optimal training duration, as described by~\citet{dubey2024llama}. While space patching offers a more straightforward solution, not requiring real-time computation of an entropy model during training, our results reveal that the improvements attributed to entropy-based patching observed in scaling experiments (Section~\ref{section:scaling}) also extend to various downstream benchmarks.\footnote{Space patching outcomes are derived from earlier experiments that did not use cross-attention; analogous patterns emerge when cross-attention is incorporated.}

\paragraph{Cross-Attention} 
In~\autoref{tab:crossattn_ablations_1b}, we perform an ablation study to evaluate the effects of inserting cross-attention at different layers of the encoder and decoder within {\textsc BLT}. For encoder cross-attention, our experiments consider three strategies for query initialization: (1) using an identical learned embedding across all global states, (2) employing a hash embedding derived from the bytes comprising the patch, and (3) applying a pooling operation to the encoder’s hidden states associated with the patch bytes at the specified encoder layer.

Our results indicate that incorporating cross-attention within the \textit{decoder} yields the highest effectiveness. Employing cross-attention in the encoder provides only marginal benefits, and these gains emerge primarily when queries are initialized using pooling. Furthermore, our analysis shows that cross-attention proves especially advantageous for the Common-Crawl dataset, with the effect being more pronounced when larger patch sizes are used.

\paragraph{n-gram Hash Embeddings} 

We conduct ablations using n-gram hash embedding vocabularies of size 0, 100K, 200K, and 400K, and summarize the outcomes in Table~\ref{tab:ngram_ablations_1b}. Our results indicate that hash embeddings contribute positively across all evaluated domains, with the most notable impact observed on Wikipedia and Github (yielding a 0.04 bpb improvement, as opposed to 0.01 bpb after 15k training steps at 8B scale). When considering the transition from 500K to 300K hashes at the 8B model size, performance differs by approximately 0.001 bpb after 15k steps. These findings suggest that the incorporation of hashes is essential for enabling {\textsc BLT} to achieve performance parity with models relying on tokenizers; nonetheless, returns become marginal beyond the 300K hash mark. Furthermore, the observed benefits of hashes seem to function independently from cross-attention, as each method boosts results in distinct datasets.

\paragraph{Local Model Hyperparameters} Table \ref{tab:local_layers} presents an ablation study on different layer counts for both the local encoder and decoder. Our findings indicate that, in conjunction with hash n-gram embeddings, {\textsc BLT} achieves optimal performance when utilizing a minimally sized encoder—specifically, a single layer—while benefiting from a comparatively deeper decoder.

\section{Related Work}
\label{section:related_work}

\textbf{Character-Level RNNs:} Character Language Modeling has long served as a foundational challenge within neural approaches to language due to its inherent capability to handle out-of-vocabulary terms in a natural way, eliminating the need for traditional back-off techniques~\citep{sutskever2011generating,mikolov2012subword,graves2013generating}. In their work, \cite{kim2016character} introduce an architecture leveraging convolutional and highway networks operating over input characters, whose outputs are then fed into LSTM-based RNNs. This model achieved parity with the then state-of-the-art RNN language models on English datasets and demonstrated superior results on languages with complex morphological structures, illustrating another notable strength of character-level LLMs. Furthering this direction, \cite{kenter2018byte} employ LSTM models at the byte level for machine comprehension, finding that such models surpass word-level counterparts particularly on morphologically complex languages such as Turkish and Russian. Likewise, character-based convolutional neural models put forth by \cite{zhang2015character} demonstrated superior outcomes in certain classification scenarios compared to word-based models. Moving toward more sophisticated architectures, \citet{chung2019hierarchical} present hierarchical LSTM models equipped with boundary detectors at each abstraction level, designed to identify latent textual hierarchies and thus enhance performance in character-level language modeling. Differently, ByteNet proposed by \cite{kalchbrenner2016neural} replaces attention with convolutional networks on character sequences for the purpose of machine translation.

\textbf{Character-Level Transformers:} The introduction of attention-based transformer architectures~\citep{vaswani2017attention}, in conjunction with subword tokenization strategies~\citep{sennrich-etal-2016-neural}, resulted in substantial advancements in neural language modeling and related benchmark tasks. Despite these improvements, the use of word or subword tokens imposes a particular inductive bias concerning the granularity at which models process linguistic input. In an effort to merge the effectiveness of transformers with earlier success in character-level modeling, \citet{al2019character} trained extremely deep transformer architectures, incorporating auxiliary loss functions to facilitate training; their transformer-based character models surpassed the performance of preceding LSTM character-level language models. Nevertheless, these models still lagged considerably behind large language models operating at the word level. Similarly, GPT-2~\citep{radfordgpt2} reported that on extensive corpora such as the 1 Billion Word Benchmark, language models trained on bytes did not attain competitiveness with those based on word-level representations.

Although \cite{Choe2019BridgingTG} established that byte-level LLMs leveraging transformer architectures can surpass the performance of subword-level LLMs with similar parameter counts, these models demand significantly higher computational resources and longer training durations. In a related vein, \cite{el2020characterbert} introduced CharFormer, a BERT-based approach in which word representations are derived through convolutional operations on character-level embeddings. Despite yielding better results within the medical domain, this method also incurs substantial computational costs. CANINE, presented by \cite{clark2022canine}, is a 150M-parameter encoder-only architecture designed to work natively with character sequences. The CANINE model incorporates a deep transformer core—akin to the global transformer used in our work—and deploys both local transformers and strided convolutions to downsample input sequences. CANINE outperforms its token-level encoder-only counterpart (mBERT) on various downstream multilingual tasks. The ByT5 model~\citep{xue2022byt5} advanced the development of byte-level encoder-decoder models without any patching strategies. Although ByT5 demonstrates increased robustness to input noise and delivers performance competitive with tokenizer-based alternatives using a quarter of the data, the absence of patching leads to computationally intensive attention operations over all bytes, substantially increasing computational demands. Modeling at the byte level, as opposed to the subword level, inevitably lengthens input sequences, presenting additional difficulties for scaling these models efficiently. Recent advances leveraging the Mamba Architecture~\citep{gu2023mamba}, which is capable of retaining a constant-size memory state across extensive contexts, have led \cite{wang2024mambabyte} to train a byte-level Mamba variant (also without patching). This approach achieves superior performance compared to byte-level transformer models under FLOP-controlled settings at the 350M parameter level, particularly in terms of bits-per-byte across several datasets.

\textbf{Patching-based approaches:} Utilizing patching mechanisms has proven effective in significantly reducing the high computational cost—measured in flops—associated with byte-level LLMs, while aiming to preserve performance quality. Early investigations, primarily focusing on smaller models and limited training bytes, demonstrated promising outcomes. For instance, \cite{nawrot-etal-2022-hierarchical} explored static patching through downsampling and upsampling within the architecture, introducing the hourglass transformer, which surpassed competing byte-level models at the 150M parameter scale. Building on these results, \cite{nawrot-etal-2023-efficient} enhanced the methodology with dynamic patching, incorporating several boundary-predictor mechanisms—one learned end-to-end, another trained with supervision from specific tokenizers, and an entropy-based model akin to {\textsc BLT}. Their experiments indicated superior performance compared to standard transformers available at that time, demonstrated on a 40M parameter model trained over approximately 400M tokens. Additionally, \cite{lester2024training} examined the impact of training on input sequences compressed via arithmetic coding—attaining compression rates that exceed those possible with BPE. Leveraging an equal-info windows method, their approach surpassed byte-level baselines in language modeling, though it fell short when compared to subword-level baselines.

This study is primarily influenced by MegaByte~\citep{yu2023megabyte}, a decoder-only causal LLM that implements a fixed, static approach to patching by concatenating representations for the transition between bytes and patches, with a local decoder-side model handling the reverse mapping from patches back to bytes. MegaByte is shown to perform on par with tokenizer-based architectures at the 1B parameter level using a dataset containing approximately 400B bytes. Throughout our experiments, we include MegaByte as a baseline and observe that its static patching mechanism falls short compared to state-of-the-art, compute-optimized tokenizer models under equivalent flop budgets; through our analysis, we demonstrate how {\textsc BLT} addresses and narrows this shortfall. Consistent findings are reported by \cite{slagle2024spacebyte}, who advocate for enhancing static patching by, for instance, patching at whitespaces or similar characters, as well as introducing a local encoder. Their modifications yield gains over tokenizer-based transformers in compute-limited settings for particular domains like Github and arXiv at the 1B parameter level. We extend this line of experimentation, highlighting that more advanced model architectures are required for byte-level models to achieve superior scaling and match the performance levels of contemporary token-based models, such as Llama 3.

\section{Limitations and Future Work}

In this study, our model training follows the recommended step count identified for Llama~3~\citep{dubey2024llama} when making architectural decisions. Yet, it should be noted that these scaling rules were originally derived for BPE-level transformers, so applying them to {\textsc BLT} may result in non-optimal ratios of data to parameter counts. Determining dedicated scaling laws tailored to {\textsc BLT}—which may reveal even more advantageous scaling patterns for this model—remains an avenue for future research. Furthermore, since the majority of our experimental results come from models with up to 1B parameters, it is plausible that ideal architectural settings will shift with scaling to 8B parameters or larger, potentially leading to enhanced outcomes at increased model sizes.

Current transformer libraries and codebases are optimized primarily for transformer architectures that rely on tokenizers. Although we include theoretical flop-matched experiments and employ some efficient solutions—like FlexAttention—for layers that differ from the standard transformer, our implementations may not achieve the same wall-clock efficiency as those based on tokenizers and could still require additional optimization to close this gap.

Although {\textsc BLT} currently relies on a separately trained entropy model for patching, an intriguing avenue for future research would be to develop patching models within an end-to-end learning framework. In Section \ref{sec:distillation}, we provide preliminary experimental results that suggest the viability of “byte-ifying” tokenizer-based architectures like Llama 3—which have been pre-trained on over 10T tokens—by transferring and freezing their transformer weights. Further exploration in this area could yield approaches that preserve the advantages conferred by bytefying, while also potentially surpassing the performance of these original tokenizer-based models, all without necessitating retraining from the ground up.

\section{Conclusion}
\label{section:conclusion}

This work introduces the Byte Latent Transformer (\textbf{BLT}), an architecture that fundamentally challenges the traditional reliance on static, fixed-vocabulary tokenization in large language models. By developing a dynamic, trainable approach for aggregating bytes into variable-length patches, {\textsc BLT} adapts its computational investment in accordance with the underlying data's complexity. This adaptive allocation brings marked gains in both computational efficiency and model resilience. Comprehensive scaling experiments indicate that {\textsc BLT} achieves performance on par with conventional tokenization-based systems such as Llama 3 when scaled up to 8B parameters and 4T bytes, while offering up to a 50\% reduction in inference flops with only minor impacts on key evaluation metrics. Additionally, {\textsc BLT} introduces a new axis for scaling, permitting concurrent increases in both model capacity and patch size without exceeding a predetermined inference budget. This capability is particularly beneficial under compute constraints typical of real-world applications. By operating directly on byte-level inputs, {\textsc BLT} also enhances the model's capacity to process rare or noisy data, leading to better robustness and improved modeling of sub-word phenomena. Collectively, these findings highlight {\textsc BLT} as a strong candidate for replacing classical tokenization-dependent models, offering a flexible, scalable, and robust paradigm for next-generation language modeling.

\end{document}